\section{Банк практических заданий}

В этом разделе собраны задачи для закрепления материала. Задания разделены по темам и уровням сложности:
\begin{itemize}
    \item \textbf{Базовый (Б):} Проверка понимания основных команд.
    \item \textbf{Продвинутый (П):} Требует написания алгоритмов.
    \item \textbf{Комплексный (К):} Мини-проект, объединяющий несколько тем.
\end{itemize}

\subsection{Блок 1: Работа с камерой и ArUco (aruco\_examples, camera\_examples)}

\begin{enumerate}
    \item \textbf{(Б) <<Фотограф>>:} Напишите скрипт, который делает фото каждые 5 секунд и сохраняет их в папку с текущей датой и временем в названии.
    \item \textbf{(Б) <<Детектор движения>>:} Сравнивайте два последовательных кадра. Если количество изменившихся пикселей превышает порог — выводите сообщение <<Движение обнаружено!>>.
    \item \textbf{(П) <<Умный светофор>>:} Распечатайте красный и зеленый круги. Если дрон видит зеленый цвет — он взлетает на 1 метр. Если красный — садится.
    \item \textbf{(П) <<Поиск маркера>>:} Дрон взлетает и вращается вокруг своей оси (yaw). Как только он видит любой ArUco маркер — останавливается и <<смотрит>> на него 3 секунды.
    \item \textbf{(П) <<Дистанция>>:} Дрон висит напротив маркера. Если маркер приближается (человек подходит с маркером в руках) — дрон отлетает назад, сохраняя дистанцию 1 метр.
    \item \textbf{(К) <<Инвентаризация>>:} В комнате разложены 5 маркеров с ID от 0 до 4. Дрон должен облететь комнату, найти все маркеры и записать в файл: <<ID 3 найден в 12:00>>, <<ID 1 найден в 12:05>>.
    \item \textbf{(К) <<Посадка на платформу>>:} Дрон должен найти маркер на полу, снизиться, зависнуть точно над центром и совершить посадку.
\end{enumerate}

\subsection{Блок 2: Телеметрия (get\_telemetry)}

\begin{enumerate}
    \item \textbf{(Б) <<Черный ящик>>:} Записывайте высоту, заряд батареи и углы наклона (roll, pitch) в CSV файл 10 раз в секунду во время полета.
    \item \textbf{(Б) <<Сигнализация>>:} Если заряд аккумулятора падает ниже 10.5 Вольт, дрон должен включить красную подсветку и начать посадку.
    \item \textbf{(П) <<Альтиметр>>:} Взлетите и висите. Выводите разницу между высотой по барометру и высотой по дальномеру (ToF). В каких ситуациях они отличаются?
    \item \textbf{(П) <<Компас>>:} Напишите программу, которая разворачивает дрон носом строго на Север (0 градусов), затем на Восток (90), Юг (180) и Запад (270).
    \item \textbf{(К) <<Виртуальная стена>>:} Задайте границы зоны полета (например, $X$ от -2 до 2 метров). Если телеметрия показывает выход за границы — дрон должен остановиться и вернуться в центр.
\end{enumerate}

\subsection{Блок 3: Миссии и полет по точкам (mission\_examples)}

\begin{enumerate}
    \item \textbf{(Б) <<Квадрат>>:} Пролетите по точкам $(0,0) \to (1,0) \to (1,1) \to (0,1) \to (0,0)$ на высоте 1 метр.
    \item \textbf{(Б) <<Лестница>>:} Пролетите вперед на 1 метр, набрав 0.5 метра высоты. Повторите 3 раза (как по ступенькам).
    \item \textbf{(П) <<Спираль>>:} Напишите генератор точек для полета по спирали (расширяющийся круг с набором высоты).
    \item \textbf{(П) <<Патруль>>:} Дрон летает между двумя точками бесконечно, пока не увидит препятствие (по датчику расстояния или камере).
    \item \textbf{(К) <<Сканирование поля>>:} Реализуйте полет <<змейкой>> над зоной $3 \times 3$ метра для создания фотоплана местности.
\end{enumerate}

\subsection{Блок 4: Групповой полет (group\_flight\_examples)}

\begin{enumerate}
    \item \textbf{(Б) <<Синхронный старт>>:} Два дрона взлетают одновременно по команде с ПК, висят 5 секунд и садятся.
    \item \textbf{(П) <<Карусель>>:} Два дрона висят напротив друг друга и меняются местами, пролетая по дуге (чтобы не столкнуться).
    \item \textbf{(П) <<Ведущий-ведомый>>:} Первый дрон летит по программе (квадрат). Второй дрон получает телеметрию первого и повторяет его движения со смещением (на 1 метр сзади).
    \item \textbf{(К) <<Шоу дронов>>:} Составьте хореографию для 4 дронов под музыку (используйте мигание светодиодами в такт).
\end{enumerate}

\subsection{Блок 5: Трекинг объектов (PioneerHumanTracking)}

\begin{enumerate}
    \item \textbf{(Б) <<Приветствие>>:} Если дрон видит лицо человека — он мигает зеленым и делает <<кивок>> (pitch вперед-назад).
    \item \textbf{(П) <<Следи за цветом>>:} Возьмите яркий оранжевый мяч. Дрон должен поворачиваться (yaw) так, чтобы мяч всегда был в центре кадра.
    \item \textbf{(П) <<Держи дистанцию>>:} Дрон видит лицо. Если лицо занимает слишком много места в кадре (человек близко) — отлетаем. Слишком мало (человек далеко) — подлетаем.
    \item \textbf{(К) <<Оператор>>:} Дрон летит боком параллельно идущему человеку, снимая его профиль (требует сложной координации движений).
\end{enumerate}

\subsection{Критерии оценки выполнения}
Задание считается выполненным, если:
\begin{enumerate}
    \item Код запускается без синтаксических ошибок.
    \item Дрон выполняет поставленную задачу в реальных условиях не менее 3 раз подряд.
    \item Код содержит комментарии, объясняющие логику работы.
    \item Соблюдены меры безопасности (дрон не вылетает за пределы зоны, экстренная посадка работает).
\end{enumerate}
